\chapter{Código fuente de la herramienta Python}
\label{anexo:python}

Este anexo recoge los fragmentos más representativos del código fuente de la
herramienta de diagnóstico implementada en Python sobre Raspberry Pi. El código
completo está disponible en el repositorio del proyecto (\texttt{src/}).

\section*{B.1 — Interfaces abstractas del núcleo}

\begin{lstlisting}[language=Python, caption={Interfaz ITransport — contrato de la capa de transporte}]
from __future__ import annotations
from abc import ABC, abstractmethod
from types import TracebackType


class ITransport(ABC):

    @abstractmethod
    def connect(self) -> None: ...

    @abstractmethod
    def disconnect(self) -> None: ...

    @abstractmethod
    def send(self, payload: bytes) -> None: ...

    @abstractmethod
    def receive(self) -> bytes: ...

    @abstractmethod
    def __enter__(self) -> ITransport: ...

    @abstractmethod
    def __exit__(
        self,
        exc_type: type[BaseException] | None,
        exc_val: BaseException | None,
        exc_tb: TracebackType | None,
    ) -> bool | None: ...
\end{lstlisting}

\section*{B.2 — Transporte ISO-TP sobre SocketCAN}

\begin{lstlisting}[language=Python, caption={IsoTpTransport — implementacion concreta sobre python-can + can-isotp}]
import time
import can
import isotp
from config.can_config import CAN_RX_ID, CAN_TX_ID, ISOTP_CF_SEPARATION_MS, ISOTP_PADDING_BYTE
from core.exceptions import DiagnosticTimeoutError
from core.interfaces.i_transport import ITransport


class IsoTpTransport(ITransport):

    def __init__(self, channel: str = "can0", timeout: float = 5.0) -> None:
        self._channel = channel
        self._timeout = timeout
        self._address = isotp.Address(
            addressing_mode=isotp.AddressingMode.Normal_11bits,
            txid=CAN_TX_ID,
            rxid=CAN_RX_ID,
        )
        self._params = {
            "stmin": ISOTP_CF_SEPARATION_MS,
            "blocksize": 0,
            "tx_padding": ISOTP_PADDING_BYTE,
        }
        self._stack: isotp.CanStack | None = None
        self._bus: can.BusABC | None = None

    def connect(self) -> None:
        self._bus = can.Bus(channel=self._channel, interface="socketcan")
        # Vaciar tramas residuales del buffer del kernel
        while self._bus.recv(timeout=0) is not None:
            pass
        self._stack = isotp.CanStack(self._bus, address=self._address,
                                     params=self._params)
        self._flush_rx()

    def disconnect(self) -> None:
        if self._stack is None:
            return
        self._bus.shutdown()
        self._stack = None
        self._bus = None

    def send(self, payload: bytes) -> None:
        self._flush_rx()
        self._stack.send(payload)
        while self._stack.transmitting():
            self._stack.process()
            time.sleep(0.001)

    def receive(self) -> bytes:
        deadline = time.monotonic() + self._timeout
        while time.monotonic() < deadline:
            self._stack.process()
            if self._stack.available():
                return self._stack.recv()
            time.sleep(0.001)
        raise DiagnosticTimeoutError(f"No response within {self._timeout:.1f}s.")

    def _flush_rx(self) -> None:
        self._stack.process()
        while self._stack.available():
            self._stack.recv()
            self._stack.process()
\end{lstlisting}

\section*{B.3 — Logger SQLite con modo WAL}

\begin{lstlisting}[language=Python, caption={SqliteDataLogger — esquema, buffer de escritura y modo WAL}]
import sqlite3, threading
from datetime import datetime, timezone
from core.interfaces.i_data_logger import IDataLogger

_SCHEMA = """
PRAGMA journal_mode=WAL;

CREATE TABLE IF NOT EXISTS sessions (
    id         INTEGER PRIMARY KEY AUTOINCREMENT,
    label      TEXT    NOT NULL DEFAULT '',
    started_at TEXT    NOT NULL,
    ended_at   TEXT
);

CREATE TABLE IF NOT EXISTS samples (
    id           INTEGER PRIMARY KEY AUTOINCREMENT,
    session_id   INTEGER NOT NULL REFERENCES sessions(id),
    pid          INTEGER NOT NULL,
    name         TEXT    NOT NULL,
    value        REAL    NOT NULL,
    unit         TEXT    NOT NULL,
    monotonic_ts REAL    NOT NULL,
    wall_ts      TEXT    NOT NULL
);

CREATE TABLE IF NOT EXISTS commands (
    id           INTEGER PRIMARY KEY AUTOINCREMENT,
    session_id   INTEGER NOT NULL REFERENCES sessions(id),
    command      TEXT    NOT NULL,
    request_hex  TEXT    NOT NULL,
    response_hex TEXT    NOT NULL,
    wall_ts      TEXT    NOT NULL
);

CREATE TABLE IF NOT EXISTS dtcs (
    id           INTEGER PRIMARY KEY AUTOINCREMENT,
    session_id   INTEGER NOT NULL REFERENCES sessions(id),
    code         TEXT    NOT NULL,
    raw_hex      TEXT    NOT NULL,
    description  TEXT    NOT NULL DEFAULT '',
    wall_ts      TEXT    NOT NULL
);
"""

_BUFFER_SIZE = 50


class SqliteDataLogger(IDataLogger):
    # WAL permite escrituras concurrentes desde hilo monitor y hilo BLE.

    def __init__(self, db_path: str = "diagnostics.db") -> None:
        self._lock = threading.Lock()
        self._sample_buffer: list[tuple] = []
        self._conn = sqlite3.connect(db_path, check_same_thread=False)
        self._conn.executescript(_SCHEMA)
        self._conn.commit()

    def log_sample(self, session_id: int, sample) -> None:
        row = (session_id, sample.pid, sample.name, sample.value,
               sample.unit, sample.timestamp,
               datetime.now(timezone.utc).isoformat())
        with self._lock:
            self._sample_buffer.append(row)
            if len(self._sample_buffer) >= _BUFFER_SIZE:
                self._flush_buffer()

    def log_command(self, session_id, command, request, response) -> None:
        with self._lock:
            self._conn.execute(
                "INSERT INTO commands"
                " (session_id, command, request_hex, response_hex, wall_ts)"
                " VALUES (?, ?, ?, ?, ?)",
                (session_id, command, request.hex(), response.hex(),
                 datetime.now(timezone.utc).isoformat()),
            )
            self._conn.commit()

    def _flush_buffer(self) -> None:
        if not self._sample_buffer:
            return
        self._conn.executemany(
            "INSERT INTO samples"
            " (session_id, pid, name, value, unit, monotonic_ts, wall_ts)"
            " VALUES (?, ?, ?, ?, ?, ?, ?)",
            self._sample_buffer,
        )
        self._conn.commit()
        self._sample_buffer.clear()
\end{lstlisting}

\section*{B.4 — Patrón decorador en dos niveles}

\texttt{LoggingTransport} envuelve el transporte y registra cada trama en el log
textual, conservando la última trama enviada/recibida por hilo. \texttt{LoggedDiagnosticSession}
envuelve la sesión y persiste en SQLite el comando ejecutado con sus tramas.

\begin{lstlisting}[language=Python, caption={LoggingTransport — registro de tramas y estado thread-local}]
import threading, logging
from core.interfaces.i_transport import ITransport
from infraestructure.logging.frame_formatter import format_rx, format_tx

_log = logging.getLogger("transport")


class LoggingTransport(ITransport):
    def __init__(self, inner: ITransport) -> None:
        self._inner = inner
        self._local = threading.local()  # last_sent/received por hilo

    @property
    def last_sent(self) -> bytes:     return getattr(self._local, "last_sent", b"")
    @property
    def last_received(self) -> bytes: return getattr(self._local, "last_received", b"")

    def send(self, payload: bytes) -> None:
        self._local.last_sent = payload
        _log.info(format_tx(payload))
        self._inner.send(payload)

    def receive(self) -> bytes:
        data = self._inner.receive()
        self._local.last_received = data
        _log.info(format_rx(data))
        return data
\end{lstlisting}

\begin{lstlisting}[language=Python, caption={LoggedDiagnosticSession — persistencia del comando en SQLite}]
class LoggedDiagnosticSession(IDiagnosticSession):
    def __init__(self, inner, logger, session_id, transport) -> None:
        self._inner = inner; self._logger = logger
        self._session_id = session_id; self._transport = transport

    def _log(self, command: str) -> None:
        self._logger.log_command(
            self._session_id, command,
            self._transport.last_sent, self._transport.last_received,
        )

    def get_engine_rpm(self) -> float:
        result = self._inner.get_engine_rpm()
        self._log("get_engine_rpm")
        return result

    def get_dtcs(self):
        result = self._inner.get_dtcs()
        self._log("get_dtcs")
        self._logger.log_dtcs(self._session_id, result)
        return result
    # ... resto de operaciones analogas ...
\end{lstlisting}

\section*{B.5 — Manejador de comandos BLE y descubrimiento de PIDs}

\begin{lstlisting}[language=Python, caption={BtCommandHandler — autenticación, despacho y probe\_pids}]
class BtCommandHandler:
    def handle(self, cmd: dict) -> dict:
        name = cmd.get("cmd", "")
        if not self._authenticated:
            if name != "auth": return {"status": "error", "message": "not authenticated"}
            if cmd.get("token") == self._auth_token:
                self._authenticated = True
                return {"status": "ok", "data": "authenticated"}
            return {"status": "error", "message": "invalid token"}

        dispatch = {
            "ping": self._ping, "snapshot": self._snapshot, "dtcs": self._dtcs,
            "clear_dtcs": self._clear_dtcs, "vin": self._vin,
            "monitor_start": self._monitor_start, "monitor_stop": self._monitor_stop,
            "sessions": self._sessions, "session_samples": self._session_samples,
            "session_commands": self._session_commands, "session_dtcs": self._session_dtcs,
            "uds_session": self._uds_session, "uds_read_did": self._uds_read_did,
            "probe_pids": self._probe_pids, "disconnect": self._cmd_disconnect,
        }
        fn = dispatch.get(name)
        if fn is None: return {"status": "error", "message": f"Unknown command: {name!r}"}
        try:    return fn(cmd)
        except Exception as exc: return {"status": "error", "message": str(exc)}

    def _probe_pids(self, _cmd: dict) -> dict:
        # 1) PIDs declarados via mascaras de soporte 0x00/0x20/.../0xE0
        declared: set[int] = set()
        for s in (0x00, 0x20, 0x40, 0x60, 0x80, 0xA0, 0xC0, 0xE0):
            with self._lock:
                self._transport.send(bytes([0x01, s]))
                raw = self._transport.receive()
            if len(raw) >= 6 and raw[0] == 0x41 and raw[1] == s:
                mask = (raw[2] << 24) | (raw[3] << 16) | (raw[4] << 8) | raw[5]
                for i in range(32):
                    if mask & (0x80000000 >> i): declared.add(s + i + 1)
                if not (mask & 0x01): break
        # 2) confirmar cada PID con una consulta real
        confirmed = {p for p in declared if self._poll_ok(p)}
        self._supported_pids = confirmed & set(PIDS.keys())
        return {"status": "ok", "data": sorted(self._supported_pids)}
\end{lstlisting}
